\documentclass[12pt,a4paper]{article}
\usepackage[utf8]{inputenc}
\usepackage[T1]{fontenc}
\usepackage{lmodern}
\usepackage{graphicx}
\usepackage{setspace}
\usepackage{geometry}

\begin{document}

\begin{center}
    \centering

    {\Large \textbf{Centro Universitário Católica de Santa Catarina}}\\[0.5cm]
    {\large Engenharia de Software}\\[4cm]

    {\Huge \textbf{Manifesto Ágil no Dia a Dia}}\\[0.5cm]

    {\large Autores:}\\[0.5cm]
    {\large Felipe Obertier Gesser}\\[0.5cm]
    {\large João Miguel de Castro Menna}\\[0.5cm]
    {\large Lucas Warmling}\\[0.5cm]
    {\large Matheus Finder Martins}\\[0.5cm]
    {\large Rafael Pavesi dos Passos}\\[2.5cm]

    {\large Professor: Edson Vaz Lopes}\\[0.5cm]

    \vfill
    {\large Joinville}\\
    {\large 2025}
\end{center}

\section{Exercício 1}

\subsection{}

\textbf{1.} O Manifesto defende "Colaboração com o cliente mais que negociação de contratos".
Como vocês imaginam que essa colaboração aconteceria idealmente em um projeto?
Que desafios podem surgir nesse processo?

\textbf{R:} Devem ocorrer reuniões frequentes de alinhamento entre ambas as partes, existindo um método
de comunicação acessível. Os desafios enfrentados podem incluir falta de alinhamento, conflito
de ideias, e dificuldade em definição de requisitos além da dificuldade na comunicação.

\subsection{}

\textbf{2.} O Manifesto Ágil encoraja equipes auto-organizadas. O que significa isso na prática?
e atitudes são necessárias para que uma equipe funcione bem nesse modelo?

\textbf{R:} Equipes auto-organizadas são aquelas que conseguem trabalhar de forma autônoma,
sem depender do controle direto de um gerente. Para isso, é essencial que tenham processos
bem definidos, alinhamento entre os colaboradores e conhecimento do fluxo de trabalho.
Habilidades como autonomia, gestão do próprio tempo, proatividade e colaboração são
indispensáveis para que o grupo funcione de maneira eficaz.

\section{Exercício 2}

\textbf{1.} Em um mercado altamente competitivo, como o equilíbrio entre a compreensão profunda
do JTBD e a busca por um TTM acelerado pode ser a chave para o sucesso de um produto, especialmente
considerando a possibilidade de lançar um produto "mínimo viável" que atenda ao JTBD central, mas
que possa ser aprimorado posteriormente com base no feedback dos usuários. Discorra sobra o tema. 

\textbf{R:} O equilíbrio entre JTBD e TTM é crucial para empresas que dependem de empréstimos
para projetos. No Brasil, o juro alto dificulta o JTBD, pois um projeto que se estende em TTM
aumenta o prejuízo de pequenas e médias empresas que não disfrutam de uma maior reserva
financeira, portanto tendo que optar por estratégias que minimizam o trabalho exaurido em
requisitos que não compõem a função principal do projeto. Logo, conciliar qualidade e tempo
é essencial para se destacar no marcado e encontrar um equilíbrio entre a compreensão profunda
do JTBD e um TTM acelerado.

\end{document}
